차분 프라이버시 수치나 워터마크 점수는 보호 대상, 실행 메커니즘, 검증
인터페이스 및 근거가 일치하지 않으면 유효한 보장이 되기 어렵다. 본 논문은 이
간극을 줄이는 주장 정합 보증(Claim-Aligned Assurance)을 제시하고, 감사 가능한
프라이버시와 안전한 공개 검증의 두 축에서 구현한다.

감사 가능한 프라이버시 연구는 완료된 윈도우 기반 DP-SGD 실행의 동결된 근거와
요청 문구를 함께 입력으로 받아, 해당 실행이 실제로 허용하는 프라이버시 단위와
수치의 범위를 판정한다. 실행 단위와 요청 단위가 일치하면 직접 경로를, 등록된
안정성 경계로 변환할 수 있으면 변환 또는 그룹 경로를 사용하며, 근거가 누락되거나
변환 결과가 무의미하면 긍정 문구를 차단한다. 통제된 100개 허용 사례와 100개
대응 비허용 사례에서 전체 판정 튜플이 예상 결과와 일치하였다. WISDM과 UCI HAR
실행은 동일한 수치라도 요청 단위에 따라 허용 범위가 달라짐을 보였고, 5,000명
규모의 Sepsis 자료와 2,779만여 건의 Backblaze 기록은 윈도우 정책에 따라 원시
단위와 생성 레코드 사이의 영향 범위가 크게 달라짐을 확인하였다. 

PP-Mark는 생성 문맥에 결속된 잠재 워터마크, 추적 커밋먼트 및 영지식 증명
영수증을 결합한다. 전체 수락은 동일한 문맥과 정확한 이미지에 대한 검출 신호와
유효한 영수증을 모두 요구한다. 10,000개 바인딩 탐색과 화이트박스 최적화는 점수
전용 공격을 드러냈지만, 평가한 위조 조건에서는 전체 수락이 관측되지 않았고
점수 모드는 변환 후 0.74--1.00의 참양성률을 유지했다.

두 연구로부터 주장 명세, 메커니즘 결속, 근거 결속형 판정 및 실패 폐쇄형 보증을
도출한다. 본 논문의 기여는 범용 인증기가 아니라, 긍정적인 프라이버시와 출처
주장을 대상, 메커니즘, 인터페이스 및 근거가 공동으로 지지하는 범위로 제한하는
검토 가능한 보증 구조에 있다.
